\section{Tabel, Media, dan Aksesibilitas}
Tabel digunakan untuk menampilkan data yang terstruktur seperti daftar mahasiswa atau produk. Gunakan \code{th} untuk header dan tambahkan atribut agar pembaca layar memahami struktur tabel. Media seperti gambar dan video perlu diberi teks alternatif agar tetap aksesibel \cite{mdnAccessibility}.

Aksesibilitas dalam HTML juga menekankan penggunaan label yang jelas dan urutan fokus yang logis. Atribut \code{alt} pada gambar membantu pengguna dengan keterbatasan visual. Dengan praktik ini, halaman web lebih ramah bagi semua pengguna.

\begin{table}[ht]
\centering
\begin{tabular}{|P{4cm}|P{7cm}|}
\hline
\textbf{Elemen} & \textbf{Catatan Aksesibilitas} \\
\hline
\code{img} & Wajib gunakan atribut \code{alt} \\
\hline
\code{table} & Gunakan \code{th} untuk header \\
\hline
\code{label} & Hubungkan ke input dengan \code{for} \\
\hline
\end{tabular}
\caption{Praktik aksesibilitas dasar HTML}
\end{table}

